\documentclass[english, parskip=full]{scrartcl}
\usepackage[english]{babel}  % Sprache auf Deutsch 
\usepackage[utf8]{inputenc}  % Kodierung der Datei
\usepackage[left=2cm, right=2cm, top=2cm, bottom=3cm, footskip=1.5cm]{geometry}
\usepackage{color}
\usepackage{graphicx, gensymb}
\usepackage{amsfonts, cancel, amssymb}
\usepackage{amsmath, amsthm, array, esint}
\usepackage{bm} % bold math symbols, e.g. \bm{\sigma} etc.
\usepackage{booktabs, multirow}  % schönere Tabellen
%\usepackage{scrlayer-scrpage}    % Manipulation des Seitenstils
%\pagestyle{scrheadings}  % Stil für die Seite setzen
\usepackage{tabularx}
\usepackage{setspace}
%\automark{section}  % Abschnittsnamen als Seitenbeschriftung 
%\setheadsepline{.5pt}    % Kopzeile durch Linie abtrennen
\usepackage[hidelinks]{hyperref}  % Links und weitere PDF-Features
\usepackage{typearea, tikz, pgffor, verbatim}
\usetikzlibrary{calc, arrows.meta,arrows, graphs, quotes, positioning, angles, babel}
\usepackage[cache=false, outputdir=build]{minted}
\usemintedstyle{vs}
\usepackage{placeins} % provides \FloatBarrier

\definecolor{bg}{rgb}{0.9,0.9,0.9}
\newcommand\numberthis{\addtocounter{equation}{1}\tag{\theequation}}
\usepackage{svg}
\usepackage{siunitx}
\usepackage{physics}
\usepackage{tensor}
\usepackage{slashed}
\usepackage{bbm}
\usepackage[compat=1.1.0]{tikz-feynman}

\usepackage{caption}
\captionsetup{
    % width=.8\textwidth,
    % indention=1cm,
    margin=0.5cm,
    format=plain,
    labelfont=bf
}
%\usepackage{subcaption}
%\subcaptionsetup{
%    indention=0cm,
%    format=hang,
%    margin=0.5cm,
%    labelfont=normalfont
%}
\usepackage[english,capitalise]{cleveref}
\usepackage[
    backend=biber,
    sorting=none,
    style=phys,
    giveninits=true,
    sortcites=true,
    citestyle=numeric-comp,
    % url=false,
    % doi=false,
    biblabel=brackets,
    % eprint=false,
    maxcitenames=3]{biblatex}
\usepackage{csquotes}
%\addbibresource{bibliography.bib}

\usepackage{mathtools}
% use \abs for absolute values
% use \abs for \left ... \right version and \abs[\bigg for \bigg version etc.
\let\abs\undefined
\let\bra\undefined
\let\braket\undefined
\DeclarePairedDelimiter\abs{\vert}{\vert}
\DeclarePairedDelimiter\kla{(}{)}
\DeclarePairedDelimiter\klb{[}{]}
\DeclarePairedDelimiter\klc{\{}{\}}
\DeclarePairedDelimiter\bra{\langle}{\rangle}
\DeclarePairedDelimiterX\brb[2]{\langle #1}{\rangle}{\delimsize\vert #2 }
\DeclarePairedDelimiterX\brc[3]{\langle #1}{\rangle}{\delimsize\vert #2 \delimsize\vert #3 }

\renewcommand{\i}{\text{i}}
\newcommand{\e}{\text{e}}
\DeclareMathOperator{\sgn}{sgn}
\newcommand{\kom}[1]{\overset{\substack{#1 \\ |}}{=}}
\newcommand{\koml}[2]{\overset{\substack{#1 \\ |}}{#2}}
\newcommand{\intx}[4]{\int_{#1}^{#2} #3 \,\mathrm{d}#4}
\newcommand{\intr}[2]{\int_{\mathbb{R}} #1 \, \mathrm{d}#2}
\newcommand{\intrnd}[3]{\int_{\mathbb{R}^{#1}} #2 \, \mathrm{d}^{#1}#3}
\newcommand{\intrpi}[2]{\int_{\mathbb{R}} #1 \, \frac{\mathrm{d}#2}{2\pi}}
\newcommand{\intrndpi}[3]{\int_{\mathbb{R}^{#1}} #2 \, \frac{\mathrm{d}^{#1}#3}{(2\pi)^{#1}}}
\newcommand{\oper}[2]{\vert #1 \rangle \langle #2 \vert}
\newcommand{\Text}[1]{\text{#1}}    % this is relevant because \text does not autocomplete to the default '\text{@}' but rather '\textbf' etc.
\newcommand{\powe}[1]{\e^{#1}}
\newcommand{\pow}[2]{{#1}^{#2}}
\newcommand{\Ket}[1]{\vert #1 \rangle}
\newcommand{\Bra}[1]{\langle #1 \vert}
\newcommand*{\Fourier}{\textsc{Fourier}\ }
\newcommand*{\F}{\mathcal{F}}
\DeclareMathOperator{\arsinh}{arsinh}
\DeclareMathOperator{\arcosh}{arcosh}
\DeclareMathOperator{\artanh}{artanh}
\DeclareMathOperator{\sinc}{sinc}
\DeclareMathOperator{\cscc}{cscc}
\DeclareMathOperator{\sinhc}{sinhc}
\DeclareMathOperator{\cschc}{cschc}
\newcommand{\conv}{\mathop{\scalebox{3.5}{\raisebox{-0.38ex}{\makebox[0.5\width][c]{$\ast$}}}}\limits}
\newcommand*{\unity}{\text{\usefont{U}{bbold}{m}{n}1}}   % operator one from :: https://tex.stackexchange.com/questions/566780/import-mathbb1-character-from-the-unicode-math-package

\renewcommand{\d}{\text{d}}
\renewcommand{\r}{\text{r}}
\renewcommand{\t}{\text{t}}
\renewcommand{\a}{\text{a}}
\renewcommand{\o}{\text{o}}
\renewcommand{\F}{\text{F}}
\renewcommand{\c}{\text{c}}
\newcommand{\A}{\text{A}}
\newcommand{\B}{\text{B}}
\newcommand{\R}{\text{R}}
\newcommand{\M}{\text{M}}
\newcommand{\m}{\text{m}}
\newcommand{\s}{\text{s}}
\newcommand{\f}{\text{f}}
\newcommand{\K}{\text{K}}
\newcommand{\hc}{\text{h.c.}}
\newcommand{\kf}{k_\F}
\newcommand{\vf}{v_\F}
\newcommand{\const}{\text{const.}}
\renewcommand{\L}{\text{L}}
\newcommand{\gray}[1]{\bgroup\color{white!40!black}#1\egroup}
\newcommand{\TODO}[1]{\bgroup\color{red!60!black}#1\egroup}
\newcommand\QnA[1]{\bgroup\color{blue}#1\egroup}
\newcommand\highlight[1]{\bgroup\color{green!60!black}#1\egroup}

\newcommand{\cabx}[3]{\begin{smallmatrix}\gray{A}&\gray{:}&#1 \\ \gray{B}&\gray{:}&#2 \\ \gray{\Xi}&\gray{:}&#3\end{smallmatrix}}
\newcommand{\zabx}[3]{\begin{smallmatrix}\gray{A}&\gray{:}&#1 \\ \gray{B}&\gray{:}&#2 \\ \gray{Z}&\gray{:}&#3\end{smallmatrix}}
\newcommand{\abx}[3]{\begin{smallmatrix}#1 \\ #2 \\ #3\end{smallmatrix}}

\newcommand{\normord}[1]{
  {:\mathrel{\mspace{1mu}#1\mspace{1mu}}:}
}
\newcommand{\varnormord}[1]{
  {\mathrel{\mspace{1mu}\substack{\ast\\[-1.5pt] \ast}#1\substack{\ast\\[-1.5pt] \ast}\mspace{1mu}}}
}

% punctuation styles (see https://tex.stackexchange.com/questions/56063/how-to-properly-typeset-footnotes-superscripts-after-punctuation-marks)
\let\origfootnote\footnote
\renewcommand{\footnote}[1]{\kern.06em\origfootnote{#1}}
\newcommand{\punctfootnote}[1]{\kern-.06em\origfootnote{#1}}

% Multiline equations
\newcommand{\bsplit}[1]{\begin{aligned}[t]#1\end{aligned}}

\newcommand*{\titel}{Integration Dictionary}
\newcommand*{\autor}{Joachim Schwardt}

\hypersetup{pdfauthor={\autor}, pdftitle={\titel}}

%\titlehead{titlehead}
\subject{Quantum Physics in 1D}
\title{\titel}
\author{\autor}
\date{\today}

\begin{document}

%\begin{titlepage}
\maketitle  % Titel setzen
%\tableofcontents  % Inhaltsverzeichnis setzen
%\end{titlepage}


\section{Gaussian integrals}
Note that $\d u^\ast\wedge\d u = \kla*{\d u_\r - \i\d u_\i} \wedge \kla*{\d u_\r + \i\d u_\i} = 2\i\d u_\r \d\u_\i$.
The following integrals can be proven by diagonalizing the matrix $A$ and transforming to the eigenbasis.
To that end we assume that $A^\dagger = A$, which also simplifies the result in the real case.
The results will also be valid for continuous results, so we choose the notation $\brc*{u}{A}{u}$ to combine both the finite and the infinite case.
Defining $\mathcal{D}u^\ast\mathcal{D}u = \prod_j \frac{\d u_j^\ast\d u_j}{2\pi\i}$ we find
\begin{align}
Z[J,J'] &= \int \powe{-\brc*{u}{A}{u} + \brb*{J}{u} + \brb*{u}{J'}} \,\mathcal{D}u^\ast\mathcal{D}u = \frac{1}{\det(A)} \powe{\brc*{J}{A^{-1}}{J'}} \label{eqn:gauss_c_jj'},\\
\bra*{u_1^\ast u_2} &= \frac{1}{Z[J,J]} \partial_{J_1^\ast} \partial_{J_2} Z[J,J] \Big\vert_{J=0} = (A^{-1})_{12} \label{eqn:gauss_c_expval_12}.
\end{align}
For us the more common integrals will instead involve real fields for which defining $\mathcal{D}u = \prod_j \frac{\d u_j}{2\pi}$ yields
\begin{align}
Z[J] &= \int \powe{-\frac{1}{2} \brc*{u}{A}{u} + \brb*{J}{u}} \,\mathcal{D}u = \frac{1}{\sqrt{\det(A)}} \powe{\frac{1}{2} \brc*{J}{A^{-1}}{J}} \label{eqn:gauss_r_j},\\
\bra*{u_1u_2} &= \frac{1}{Z[J]} \partial_{J_1} \partial_{J_2} Z[J] \Big\vert_{J=0} = (A^{-1})_{12} \label{eqn:gauss_r_expval_12}.
\end{align}
\section{Path integrals}
For the typical quadratic Luttinger Liquid Hamiltonian
\begin{align}
H_0 &= \frac{u}{2\pi} \intr{\klb*{K\kla*{\nabla\theta}^2 + \frac{1}{K} \kla*{\nabla\phi}^2} }{x}
\end{align}
one finds after a somewhat involved calculation that
\begin{align}
\frac{1}{K} \bra*{\klb*{\phi(r) - \phi(0)}^2} &= K \bra*{\klb*{\theta(r) - \theta(0)}^2} =  F_1(r),\\
\bra*{\phi(r)\theta(0)} &= \frac{1}{2} F_2(r)
\end{align}
where $r=(u\tau,x)$ and at finite temperature and to lowest order in $\alpha$ with $y_\alpha \equiv u\tau + \alpha\sgn(\tau)$
\begin{align}
F_1(r) &= \frac{1}{2} \log \klb*{\kla*{\frac{\beta u}{\pi\alpha}}^2 \kla*{\sin^2\kla*{\frac{\pi y_\alpha}{\beta u}} + \sinh^2\kla*{\frac{\pi x}{\beta u}} }}, \label{eqn:F1r}\\
F_2(r) &= -\i\arg\klb*{\tan\kla*{\frac{\pi y_\alpha}{\beta u}} + \i\tanh \kla*{\frac{\pi x}{\beta u}}}.\label{eqn:F2r}
\end{align}
For $T=0$ these expression can be simplified to
\begin{align}
F_1(r) &= \log \klb*{\frac{\sqrt{y_\alpha^2 + x^2}}{\alpha}},\label{eqn:F1r_T0}\\
F_2(r) &= -\i\arg \klb*{y_\alpha + \i x} \label{eqn:F2r_T0}.
\end{align}
This corresponds to the real and imaginary part of $\log(z)$ with $z=\frac{y_\alpha - \i x}{\alpha}$.
\subsection{General expectation values}
One can show that for sequences $A_j,B_j$
\begin{align}
\bra*{\powe{\i \sum_j \klb*{A_j\phi(r_j) + B_j\theta(r_j)}}} &= \powe{\frac{1}{2} \sum_{j<k} \klb*{\kla*{KA_jA_k + \frac{1}{K}B_jB_k} F_1(r_j-r_k) - (A_jB_k + A_kB_j) F_2(r_j-r_k)}} \label{eqn:expval_Aphi_Btheta_sum}.
\end{align}
This only holds assuming the neutrality condition $\sum_jA_j = 0 = \sum_jB_j$, otherwise the expectation value works out to zero.





%\begin{thebibliography}{9}
%\bibitem{example} description, \url{https://www.exmapleurl}, day.\,month~2021.
%\end{thebibliography}

\end{document}